%% Virasoro residues and their motivic comparison problem.

\providecommand{\MZV}{\mathrm{MZV}}
\providecommand{\mot}{\mathrm{mot}}
\providecommand{\Vir}{\mathrm{Vir}}
\providecommand{\Walg}{\mathcal W}
\providecommand{\cA}{\mathcal A}
\providecommand{\ClaimStatusDefinitional}{}
\providecommand{\ClaimStatusConditional}{}
\providecommand{\ClaimStatusOpen}{}

\chapter{Virasoro Residues and Motivic Comparison}
\label{ch:virasoro-motivic-purity}

\section*{Local algebra, residues, and periods}

The Virasoro local layer is exact.  The operator product, the sphere
Ward recursion, and the level-four vacuum Gram form give
\[
 T(z)T(w)\sim \frac{c/2}{(z-w)^4}
 +\frac{2T(w)}{(z-w)^2}+\frac{\partial T(w)}{z-w},
 \qquad
 N_\Lambda(c)=\frac{c(5c+22)}{10}.
\]
Consequently every connected stress-tensor Ward function is a
rational function on an ordered configuration space.  A scalar
shadow coordinate requires the additional ordered-residue package
\(\mathsf H_{\mathrm{res}}(\Vir_c;X)\).  A motivic statement requires
another comparison from the residue complex to a specified motivic
period complex.  Rationality of the local integrand and Tate
factorization of its period are distinct assertions.

\section{Regularity and the comparison packages}
\label{sec:vmp-setup}

\begin{definition}[Level-four and residue-regular charges]
\ClaimStatusDefinitional
\label{def:vmp-generic-c}
The level-four regular locus is
\[
 C^{(4)}_{\mathrm{reg}}
 :=\{c\in\mathbb C: c(5c+22)\ne0\}.
\]
For an arity bound \(R\), a charge is
\emph{residue-regular through arity \(R\)} when the package
\(\mathsf H_{\mathrm{res}}^{\le R}(\Vir_c;X)\) exists and every
scalar contraction used through arity~\(R\) is defined.  The
all-arity locus \(C^{\mathrm{res}}_{\mathrm{reg}}\) consists of
charges carrying compatible packages for every~\(R\).
Once such a compatible family is fixed, write
\(C^{\mathrm{gen}}:=C^{\mathrm{res}}_{\mathrm{reg}}\) for the
all-arity domain used by downstream comparison statements.
\end{definition}

\begin{remark}[Radical and quotient lanes]
\label{rem:vmp-exclusions}
At \(c=0\) and \(c=-22/5\), the level-four invariant form has a
radical.  The corresponding simple quotient carries its own residue
package, built after passage to the radical quotient.  Higher Kac
degeneracy loci are treated by the same typed quotient construction.
\end{remark}

\begin{problem}[Compare the formal Riccati lane with geometric residues]
\ClaimStatusOpen
\label{thm:vmp-riccati-vir-recap}
\emph{Type signature:}
\textup{(Open quadrant, formal weighted-series and ordered residue-bar
presentations, level \(2\),
\(\mathsf H_{\mathrm{res}}+\mathsf H_{\mathrm{Ricc}\to\mathrm{res}}\)).}

A formal square-root series may be generated from chosen Virasoro
seeds.  Construct a chain map from its coefficient model to the
transferred scalar residue model and prove compatibility with the
Maurer--Cartan equations, arity conventions, and normalizations.  The
package \(\mathsf H_{\mathrm{Ricc}\to\mathrm{res}}\) records this
comparison.
\end{problem}

\begin{definition}[Motivic residue package]
\ClaimStatusDefinitional
\label{def:vmp-motivic-lift}
For a constructed ordered-residue package, the motivic comparison
package \(\mathsf H_{\mathrm{mot}}\) consists of:
\begin{enumerate}[(i)]
\item a graded motivic period complex \(\mathcal P^{\mot}\) with its
de Rham coaction;
\item a collision-compatible chain map from the ordered correlator
complex to \(\mathcal P^{\mot}\);
\item a comparison from \(\mathcal P^{\mot}\) to the scalar
contraction in \(\mathsf H_{\mathrm{res}}\);
\item compatibility with the ordered-to-symmetric averaging map and
with the period homomorphism;
\item declared Tate twists, weight conventions, and regularizations.
\end{enumerate}
The resulting class is denoted
\[
 S_r^{\mot}(\Vir_c;\mathsf H_{\mathrm{res}},\mathsf H_{\mathrm{mot}}).
\]
\end{definition}

\section{The conditional purity criterion}
\label{sec:vmp-main}

\begin{theorem}[Tate-factorization criterion for Virasoro residues]
\ClaimStatusConditional
\label{thm:virasoro-motivic-purity}
\emph{Type signature:}
\textup{(Open/motivic quadrants, ordered residue and motivic-period
presentations, levels \(2\leftrightarrow5\),
\(\mathsf H_{\mathrm{res}}+\mathsf H_{\mathrm{mot}}
+\mathsf H_{\mathrm{Tate}}\)).}

Assume \(\mathsf H_{\mathrm{res}}\) and
\(\mathsf H_{\mathrm{mot}}\).  Suppose the Tate-factorization package
\(\mathsf H_{\mathrm{Tate}}\) supplies, for every~\(r\), a
factorization of the motivic residue class through the unit-motive
summand
\[
 \mathbb Q(0)\otimes\mathbb Q(c)
 \longrightarrow \mathcal P^{\mot}\otimes\mathbb Q(c).
\]
Then
\[
 S_r^{\mot}(\Vir_c)
 =S_r(\Vir_c;\mathsf H_{\mathrm{res}})\otimes1
 \in\mathbb Q(0)\otimes\mathbb Q(c),
\]
and its reduced positive-weight motivic coaction vanishes.
\end{theorem}

\begin{proof}
The factorization places the class in the unit-motive summand.  The
period comparison identifies its scalar coefficient with the
ordered-residue contraction.  The reduced positive-weight coaction
vanishes on \(\mathbb Q(0)\), which proves the assertion.
\end{proof}

\begin{remark}[Formal algebraicity and motivic periods]
\label{rem:vmp-convergence-of-sqrt}
Formal algebraicity controls coefficients of the auxiliary formal
series.  The package \(\mathsf H_{\mathrm{mot}}\) controls periods of
the geometric residue cycles.  The comparison
\(\mathsf H_{\mathrm{Ricc}\to\mathrm{res}}\) connects these two
statements.
\end{remark}

\section{Denominator and growth obligations}
\label{sec:vmp-denominator}

\begin{problem}[Determine the divisor of the residue coefficients]
\ClaimStatusOpen
\label{prop:denominator-structure}
\emph{Type signature:}
\textup{(Open quadrant, scalar residue presentation, level \(2\),
\(\mathsf H_{\mathrm{res}}+\mathsf H_{\mathrm{den}}\)).}

Construct the scalar coordinates and determine their reduced divisors
as rational functions of~\(c\).  The denominator package
\(\mathsf H_{\mathrm{den}}\) contains a proof of rational dependence,
an exact divisor calculation in each arity, and a uniform induction
compatible with the transferred brackets.  The local factor
\(c(5c+22)\) enters through the norm of~\(\Lambda\); its
multiplicity in arity~\(r\) is part of this problem.
\end{problem}

\begin{remark}[Uniform asymptotic package]
\label{rem:vmp-growth}
Growth in~\(r\), Laurent order at \(c=\infty\), and the analytic
radius of a generating series are governed by
\(\mathsf H_{\mathrm{asymp}}\).  This package includes uniform bounds
and an order-of-limits theorem.
\end{remark}

\begin{remark}[The free-boson specialization]
\label{rem:vmp-c1-excluded}
At \(c=1\), the exact local Gram factor specializes to
\(N_\Lambda(1)=27/10\).  Higher residue values at this charge are
obtained by specializing a constructed residue tower.  A comparison
with categorical coefficients in Volume~III requires a named
cross-volume chain map.
\end{remark}

\section{Rational transfer as a conditional algebraic statement}
\label{sec:vmp-transport}

\begin{theorem}[Rationality under a finite scalar transfer package]
\ClaimStatusConditional
\label{thm:virasoro-riccati-transport-rationality}
\emph{Type signature:}
\textup{(Open quadrant, transferred scalar presentation, level \(2\),
\(\mathsf H_{\mathrm{res}}+\mathsf H_{\mathrm{rat}}\)).}

Assume the transferred scalar Maurer--Cartan equation is triangular in
arity, every transfer kernel lies in \(\mathbb Q(c)\), and every
linear coefficient solved at the next arity is invertible in
\(\mathbb Q(c)\).  Then all recursively constructed scalar
coordinates lie in \(\mathbb Q(c)\).
\end{theorem}

\begin{proof}
The initial coordinates belong to \(\mathbb Q(c)\) by hypothesis.
At arity~\(r\), triangularity expresses the next coordinate as the
inverse linear coefficient times a finite polynomial in earlier
coordinates and transfer kernels.  The field \(\mathbb Q(c)\) is
closed under these operations.  Induction on~\(r\) gives the result.
\end{proof}

\begin{remark}[Local binary OPE and transferred brackets]
\label{rem:vmp-binary-bracket-closure}
The Virasoro OPE determines the local binary collision terms.  Higher
transferred brackets encode homotopies of the residue complex.  The
package \(\mathsf H_{\mathrm{rat}}\) proves their scalar kernels and
thereby turns local rationality into rationality of the transferred
coordinates.
\end{remark}

\section{Motivic coaction}
\label{sec:vmp-no-mzv}

\begin{corollary}[Unit-motive coaction]
\ClaimStatusConditional
\label{cor:virasoro-no-mzv-contribution}
\emph{Type signature:}
\textup{(Open/motivic quadrants, motivic scalar presentation,
levels \(2\leftrightarrow5\),
\(\mathsf H_{\mathrm{res}}+\mathsf H_{\mathrm{mot}}
+\mathsf H_{\mathrm{Tate}}\)).}
Under the hypotheses of
Theorem~\ref{thm:virasoro-motivic-purity}, every Virasoro motivic
residue has zero reduced positive-weight coaction.
\end{corollary}

\begin{proof}
Theorem~\ref{thm:virasoro-motivic-purity} places every class in
\(\mathbb Q(0)\otimes\mathbb Q(c)\), and the reduced
positive-weight coaction vanishes on this summand.
\end{proof}

\begin{remark}[Grothendieck--Teichm\"uller action]
\label{rem:vmp-grt-action-trivial-on-vir}
Under \(\mathsf H_{\mathrm{Tate}}\), the motivic Galois action on the
Virasoro scalar residue factors through the action on the unit
motive.  Constructing \(\mathsf H_{\mathrm{Tate}}\) is the geometric
content of this statement.
\end{remark}

\section{Comparison with principal \texorpdfstring{$\mathcal W$}{W}-algebras}
\label{sec:vmp-w-algebra-comparison}

\begin{remark}[The first additional generator lane]
\label{rem:vir-vs-w3-contrast}
Virasoro has the stress-tensor generator and rational Ward
correlators.  Principal \(\mathcal W_3\) adds a spin-three generator
and a multivariable residue complex.  Each generator lane acquires a
motivic meaning through its own collision-compatible period map and
averaging comparison.
\end{remark}

\begin{remark}[Finite-lane comparison for \texorpdfstring{$\mathcal W_N$}{WN}]
\label{rem:vmp-wn-propagation-bound}
Principal \(\mathcal W_N\) has generator weights
\(2,3,\ldots,N\).  The stress-tensor lane receives the Virasoro local
input, while every additional lane contributes separate transfer
kernels and period cycles.  The multivariable package records their
mixed brackets.
\end{remark}

\begin{problem}[Transport the motivic package through Drinfeld--Sokolov reduction]
\ClaimStatusOpen
\label{cor:vir-purity-propagates-to-dress-reduced}
\emph{Type signature:}
\textup{(Open/motivic quadrants, Drinfeld--Sokolov and ordered-residue
presentations, levels \(1\leftrightarrow2\leftrightarrow5\),
\(\mathsf H_{\mathrm{DS}\to\mathrm{res}}
+\mathsf H_{\mathrm{mot}}\)).}
Construct a chain map from the Drinfeld--Sokolov complex to the
ordered-residue complexes of the source and reduced
\(\mathcal W_N\)-algebra.  Prove compatibility with scalar
contractions, motivic period maps, and the stress-tensor projection.
This comparison determines the image of the Virasoro Tate package
under reduction.
\end{problem}

\section{The comparison frontier}
\label{sec:vmp-epilogue}

The exact Ward and Gram calculations provide the local source.  The
ordered-residue package constructs the scalar invariant.  The motivic
package lifts its chain representative to a period class.  Tate
factorization, denominator control, and Drinfeld--Sokolov propagation
are the three named obligations that complete Virasoro motivic
purity.
